\chapter{Referencial Teórico}

Este capítulo reúne os fundamentos conceituais que sustentam a proposta do chatbot financeiro, contemplando soluções correlatas e as principais tecnologias empregadas na construção do sistema.

\section{Trabalhos Relacionados}

O levantamento de trabalhos relacionados tem como objetivo identificar abordagens já utilizadas em soluções de automação conversacional, controle financeiro digital e integração entre assistentes virtuais e plataformas de mensageria. Essa análise permite compreender práticas adotadas no mercado e na literatura, bem como delimitar o diferencial da proposta desenvolvida neste trabalho.

\subsection{Concorrentes}

Entre as soluções concorrentes, destacam-se aplicativos de gestão financeira que oferecem categorização de despesas, relatórios e metas, além de assistentes conversacionais voltados a atendimento automatizado. De modo geral, essas ferramentas evidenciam a importância da usabilidade, da integração com canais familiares e da apresentação clara de indicadores financeiros. Como texto-base provisório, esta subseção deverá ser refinada com a descrição comparativa de produtos específicos e suas funcionalidades.

\subsection{Artigos Científicos Relacionados}

Na literatura científica, estudos sobre chatbots, interfaces conversacionais e apoio à tomada de decisão mostram que sistemas baseados em linguagem natural podem reduzir barreiras de uso e aumentar a frequência de interação dos usuários com serviços digitais. Para a versão final do trabalho, esta subseção deverá incorporar referências bibliográficas que abordem chatbots financeiros, interação humano-computador, adoção de mensageria e recomendação personalizada em contexto financeiro.

\section{Tecnologias}

As tecnologias selecionadas para a solução foram escolhidas com base em integração, facilidade de automação, escalabilidade e aderência ao fluxo proposto para o produto.

\subsection{WhatsApp}

O WhatsApp será utilizado como principal canal de interação com o usuário. A escolha dessa plataforma se justifica por sua ampla adoção, familiaridade de uso e capacidade de reduzir a fricção no registro de informações do dia a dia. No contexto da solução, o aplicativo funciona como porta de entrada para mensagens, comandos, alertas e confirmações de lançamentos financeiros.

\subsection{Evolution API}

A Evolution API atua como camada de integração entre o sistema proposto e o WhatsApp, permitindo o envio e recebimento de mensagens, além da orquestração de eventos conversacionais. Nesta etapa do trabalho, a subseção registra seu papel arquitetural e deverá ser expandida posteriormente com detalhes de autenticação, webhooks e fluxo de mensagens.

\subsection{Chatbot}

O chatbot representa o núcleo conversacional da solução, sendo responsável por interpretar mensagens em linguagem natural, direcionar o fluxo de atendimento e acionar os serviços internos necessários ao registro e à consulta de dados financeiros. Sua modelagem considera clareza de resposta, simplicidade operacional e contextualização adequada para o usuário final.

\subsection{n8n}

O n8n será empregado como ferramenta de automação para integrar etapas do fluxo conversacional, regras de negócio e serviços externos. Seu uso favorece a criação de rotinas visuais de processamento, reduzindo o acoplamento direto entre todas as integrações e permitindo maior flexibilidade na manutenção dos fluxos.

\subsection{Groq API}

A Groq API é considerada na proposta como serviço de apoio ao processamento de linguagem natural, com potencial para interpretar intenções, estruturar informações extraídas das mensagens e apoiar respostas contextualizadas. Esta subseção permanece como base conceitual e deverá ser aprofundada na versão final com a estratégia de uso da API na arquitetura.

\subsection{Docker}

O Docker será utilizado para padronizar a execução dos serviços da aplicação em ambientes isolados, favorecendo reprodutibilidade e portabilidade. Seu uso contribui para simplificar a configuração do ambiente de desenvolvimento e implantação da solução.

\subsection{Docker Compose}

O Docker Compose será empregado para orquestrar os contêineres necessários ao funcionamento do sistema, como backend, automação, banco de dados e interface web. A adoção dessa ferramenta permite definir e subir a infraestrutura local de forma integrada.

\subsection{React}

O React será utilizado na construção da interface web destinada à visualização de relatórios, gráficos e dados financeiros consolidados. A escolha se deve à flexibilidade para composição de componentes e à facilidade de criação de interfaces dinâmicas voltadas à experiência do usuário.

\subsection{PostgreSQL}

O PostgreSQL será o sistema de gerenciamento de banco de dados da aplicação, responsável por persistir usuários, lançamentos financeiros, categorias, metas e demais entidades necessárias ao funcionamento da solução. Sua adoção se justifica pela robustez, confiabilidade e aderência a aplicações transacionais.
